\documentclass[stu,floatsintext]{apa7}
\usepackage[american]{babel}
\usepackage{csquotes}
\usepackage[style=apa,sortcites=true,sorting=nyt,backend=biber]{biblatex}
\addbibresource{references.bib}

\title{Title of the Paper}
\shorttitle{Short Title}

\authorsnames{Student Name}
\authorsaffiliations{{Department of Psychology, University}}

\course{PSY 101: Introduction to Psychology}
\professor{Professor Name}
\duedate{Month Day, 2026}

\abstract{State the research question, the participants, the method, the main
findings with effect sizes, and the conclusions in up to 250 words.}

\keywords{first keyword, second keyword, third keyword}

\begin{document}

\maketitle

Introduce the problem, review the relevant literature and state the
hypotheses. Cite in parentheses \parencite{cohen1988statistical} or in the
text, as \textcite{opensciencecollaboration2015estimating} found. The
manuscript follows the seventh edition of the APA manual
\parencite{apa2020publication}.

\section{Method}
\subsection{Participants}
Describe the sample and how it was recruited.

\subsection{Materials and Procedure}
Describe the measures and what participants did.

\section{Results}
Table~\ref{tab:results} reports the descriptive statistics.

\begin{table}[htbp]
  \caption{Descriptive Statistics by Condition}
  \label{tab:results}
  \begin{tabular}{lrr}
    \toprule
    Condition & \textit{M} & \textit{SD} \\
    \midrule
    Control & 4.21 & 1.02 \\
    Treatment & 4.87 & 0.96 \\
    \bottomrule
  \end{tabular}
\end{table}

\section{Discussion}
Interpret the results and state the limitations.

\printbibliography

\end{document}
